\documentclass[UTF8,a4paper,12pt,fontset=fandol]{ctexart}
\usepackage{geometry}
\usepackage{enumitem}
\usepackage{setspace}
\usepackage{hyperref}

\geometry{left=2.6cm,right=2.6cm,top=2.4cm,bottom=2.4cm}
\setlength{\parindent}{0pt}
\setlength{\parskip}{0.7em}
\setstretch{1.25}
\hypersetup{hidelinks}
\pagestyle{empty}

\begin{document}
\thispagestyle{empty}

\begin{enumerate}[leftmargin=2em,itemsep=1.1em]
\item 图表清晰度不够，部分图片和公式显示不清。

答：后续会统一重新导出高清图片，检查图中文字、线宽、标注和分辨率，保证论文打印版和电子版中都能清楚辨认。

\item 公式含义说明不足，尤其是公式 2-4 等位置。

答：会在公式前后补充变量定义、符号含义和公式作用说明。对于关键公式，会进一步说明它在模型训练或实验分析中的具体作用，避免只给出公式而缺少解释。

\item 第一章缺少整体研究框图或流程图。

答：会在第一章增加论文整体研究路线图，说明研究背景、核心方法、实验验证和系统实现之间的关系，使全文结构更加清晰。

\item DeepSeek-V3 实验中内存占用较大、精度提升不明显，需要解释原因。

答：会补充该实验结果的分析。主要从模型参数规模、推理资源消耗、任务适配性和数据特征几个方面解释：大模型在通用语义理解上有优势，但在本文这种特定图像重建或识别任务中，并不一定能直接带来明显精度提升。同时会说明该结果对本文方法选择的意义，即本文更关注面向具体任务的轻量化和针对性优化，而不是单纯扩大模型规模。

\item 参考文献中 arXiv 文献较多，近三年正式发表文献偏少。

答：会重新筛选参考文献，减少预印本文献比例，补充近三年正式发表的期刊、会议论文和相关代表性工作。

\item 图像重建与语义分割之间的关系不够明确。

答：会在方法部分补充二者的逻辑关系。图像重建主要用于提升退化图像的清晰度和结构完整性，为后续语义分割提供更稳定的输入；语义分割则用于验证重建结果是否真正改善了目标区域的识别效果。后续会通过流程图和实验结果进一步体现二者之间的联系。

\item 对比方法偏旧，实验说服力还可以增强。

答：会补充近年的相关对比方法，并重新整理实验表格，使本文方法与现有方法之间的性能差异更加直观。

\item 损失函数包含多项约束，如何权衡需要说明。

答：会补充损失函数权重设置依据。本文会说明各损失项分别约束图像结构、细节信息和整体一致性，并解释为什么不能简单地继续增加损失项。损失项并不是越多越好，关键是每一项都要对应明确的优化目标，否则可能造成训练目标冲突或模型收敛不稳定。

\item 页面留白较多，部分图表排版不够紧凑。

答：会统一调整图表大小、段落间距和浮动位置，减少大面积空白，使版面更加协调。

\item 目录层级较碎，小标题过多。

答：会合并部分小节标题，减少过细的层级划分，使章节主线更加集中。

\item 缩略语存在重复定义。

答：会统一检查全文缩略语，保证首次出现时给出完整定义，后文不再重复解释。

\item 表 4-1 表达方式可以优化。

答：会重新检查表 4-1 的内容和展示方式。如果表格表达不够直观，会考虑改为图表结合的形式，使结果更容易理解。

\item 第二章中基础网络介绍过多，部分图可以删减。

答：会压缩基础网络介绍，删除与本文创新点关系不大的基础结构图，把篇幅更多放在本文方法设计和改进部分。

\item CUDA、GPU 与实验性能之间的关系需要说清楚。

答：会补充 CUDA 和 GPU 在本文实验中的具体作用，说明其主要用于加速模型训练和推理过程。后续会增加 GPU 与 CPU 的耗时对比数据，用定量结果说明 GPU 加速带来的效率提升，而不是只停留在文字描述。

\item 散斑鲁棒性实验依据不足。

答：会补充散斑噪声的设置依据、叠加方式和实验参数，并说明该实验用于验证模型在复杂噪声条件下的稳定性。后续会增加不同散斑强度下的实验结果，使鲁棒性分析更加充分。

\item 可视化结果差异不明显，图片偏小。

答：会放大关键可视化结果图，并增加局部细节对比区域。对于差异不明显的结果，会通过局部放大、误差图或指标对比辅助说明，使不同方法之间的效果差异更容易看出来。
\end{enumerate}

\end{document}
